\chapter{체킹과 전체 탐색: \texorpdfstring{\code{check}, \code{search}}{check, search}}\label{ch:rust-check}

봄베의 정지 하나가 알려 주는 것은 로터 순서, 드럼 위치, 테스트 글자의 스테커 후보뿐이다. 그리고 정지의
대부분은 거짓 정지다(\ref{sec:f-register}절). \code{check} 모듈은 정지 하나를 받아 거짓이면 버리고, 참이면 키 전체와 평문을 내놓는다.
실제로는 체킹 머신이라는 별도의 기계와 사람의 손이 하던 일이다. \code{search} 모듈은 이 모든 것을 여러 로터
순서에 대해 병렬로 돌린다.

\section{네 단계}

\code{check\_stop}은 네 단계로 이루어진다(그림~\ref{fig:g-check}).

\begin{figure}[ht]
\centering
\begin{tikzpicture}[x=1cm,y=1cm]
  \node[blkr,minimum width=18mm,minimum height=13mm] (s0) at (0,0) {정지};
  \node[blkb,minimum width=27mm,minimum height=13mm] (s1) at (3.1,0) {1. 스테커 가설\\{\footnotesize 전류를 흘려 읽기}};
  \node[blkb,minimum width=27mm,minimum height=13mm] (s2) at (6.5,0) {2. 크립 전체\\{\footnotesize 대조하며 채우기}};
  \node[blkg,minimum width=27mm,minimum height=13mm] (s3) at (9.9,0) {3. 링 설정\\{\footnotesize 넘어감 시점}};
  \node[blkg,minimum width=27mm,minimum height=13mm] (s4) at (13.3,0) {4. 남은 스테커\\{\footnotesize 언어 통계}};
  \draw[->] (s0) -- (s1); \draw[->] (s1) -- (s2); \draw[->] (s2) -- (s3); \draw[->] (s3) -- (s4);
  \draw[decorate,decoration={brace,amplitude=5pt,mirror}] (1.8,-0.95) -- (7.8,-0.95) node[midway,below=6pt,note] {논리로 확정 (기각 가능)};
  \draw[decorate,decoration={brace,amplitude=5pt,mirror}] (8.6,-0.95) -- (14.6,-0.95) node[midway,below=6pt,note] {점수로 고름};
\end{tikzpicture}
\caption{체킹의 네 단계.}
\label{fig:g-check}
\end{figure}

\section{1단계: 가설 하나에 전류를 흘린다}

\begin{lstlisting}[language=Rust,caption={\code{check.rs}: 스테커 가설(일부).}]
for v in 0..N as u8 {
    let live = plugged.trace(&stop.wheel_order, stop.position, v, &opts);
    if live.iter().any(|m| m.count_ones() > 1) {
        continue;
    }
    let mut map = [UNKNOWN; N];
    let mut ok = true;
    for (a, &m) in live.iter().enumerate() {
        if m != 0 {
            ok &= set_pair(&mut map, a as u8, m.trailing_zeros() as u8);
        }
    }
    if ok {
        out.push(map);
    }
}
\end{lstlisting}

테스트 글자의 스테커 짝 후보 26가지(\code{v}) 각각에 대해 그 선 하나에 전압을 걸고, 조기 종료 없이 끝까지
퍼뜨린다(\code{trace}). 그리고 케이블마다 켜진 선의 수를 센다(\code{count\_ones}).

\Why 참인 가설에서는 참인 귀결만 나온다. 글자 하나의 스테커 짝은 하나뿐이므로, 가설이 참이면 어떤 케이블에도
선이 \emph{둘 이상} 켜질 수 없다. 둘 이상 켜진 케이블이 하나라도 있으면 그 가설은 거짓이다(3--5행).
살아남은 가설에서는 케이블마다 켜진 선이 그 글자의 스테커 짝이다. \code{m.trailing\_zeros()}는 ``켜진 비트의
위치''를 읽는 것이다. \code{set\_pair}는 $a$의 짝이 $b$이면 $b$의 짝도 $a$라고 표에 적되, 이미 적힌 값과
어긋나면 실패를 알린다.

\General \ref{sec:f-register}절의 스도쿠와 같다. 가정에서 출발해 귀결을 끝까지 밀어 보고, 같은 줄에 같은 수가
둘 생기면 가정을 버린다.

\Ours 이것은 봄베가 하는 일과 방향만 다르다. 봄베는 틀린 가설 하나로 나머지를 한꺼번에 기각했고, 체킹은
남은 후보에 전압을 걸어 그 귀결을 \emph{읽는다}. 실제 체킹 머신에서는 조작자가 스테커 후보의 키를 누르고
켜진 램프를 읽어 메뉴의 선을 하나씩 따라갔다~\cite[12--16]{us6812report1944}. 여기서는 그것을
\code{propagate}가 한 번에 한다.

\section{2단계: 크립 전체와 대조한다}

메뉴에는 크립의 일부만 들어 있다. 스크램블러 수가 모자라 뺀 선도 있고 떨어진 덩어리의 선도 있다.
\code{extend\_with\_crib}은 크립의 \emph{모든} 글자에 대해 에니그마의 식
$\Stk(c) = \Scr_k(\Stk(p))$를 확인한다.

\begin{lstlisting}[language=Rust,caption={\code{check.rs}: 크립 한 글자의 대조(일부).}]
let (sp, sc) = (map[pl as usize], map[ci as usize]);
if sp != UNKNOWN {
    let want = s[sp as usize];
    if sc == UNKNOWN {
        if !set_pair(&mut map, ci, want) { return None; }
        changed = true;
    } else if sc != want {
        return None;
    }
} else if sc != UNKNOWN {
    if !set_pair(&mut map, pl, s[sc as usize]) { return None; }
    changed = true;
}
\end{lstlisting}

평문 글자의 짝 \code{sp}를 알면 암호문 글자의 짝은 스크램블러 표 \code{s}에서 계산된다(\code{want}).
암호문 글자의 짝을 아직 모르면 채워 넣고, 이미 알고 있는데 다르면 모순이므로 이 정지를 버린다
(\code{return None}). 반대 방향도 같다. 이 일을 새로 채워지는 것이 없을 때까지 되풀이한다. 하나를 채우면
그것이 다른 글자를 채우는 실마리가 되기 때문이다.

\Ours 거짓 정지의 대부분이 이 단계에서 탈락한다. 예제에서 157개의 정지 가운데 통과하는 것은 하나다.

\section{3단계: 링 설정}

봄베가 찾은 것은 코어 위치(창 글자 $-$ 링 설정)다. 창 글자와 링 설정을 따로 알려면 정보가 더 필요하다.

\Why 링 설정이 암호화에 영향을 주는 길은 하나뿐이다. 노치가 링에 붙어 있으므로, 링 설정은 로터가 옆 로터를
\emph{언제} 넘기는지를 정한다. 따라서 링 설정은 넘어감이 일어나는 긴 전문을 복호화해 보아야 알 수 있다.
오른쪽 로터의 링은 가운데 로터가 언제 움직이는지를 정하고, 가운데 로터의 링은 왼쪽 로터가 언제 움직이는지를
정한다. 왼쪽 로터의 링은 아무것도 넘기지 않으므로 어떤 값이든 결과가 같다. 원리적으로 알아낼 수 없고,
알 필요도 없다.

코드는 오른쪽과 가운데 링의 $26 \times 26$가지 조합을 모두 시험한다. 조합마다

\begin{enumerate}
\item 크립 구간에서 로터가 움직이는 모양이 봄베의 가정과 맞는지 확인하고,
\item 크립 위치에서 전문의 처음으로 로터를 되감아 시작 위치를 구하고(\code{wind\_back}),
\item 전문 전체를 복호화해 크립 자리에 크립이 나오는지 확인하고,
\item 복호문이 독일어다운 정도를 점수로 매긴다.
\end{enumerate}

점수가 가장 높은 조합을 택한다. 정해지지 않는 링은 \texttt{A}로 둔다.

\Ours 예제의 실제 키는 링 \texttt{GMC}, 시작 위치 \texttt{KDC}이고 복원된 키는 링 \texttt{AAC}, 시작 위치
\texttt{ERC}다. 글자는 다르지만 틀린 것이 아니다. 창 글자에서 링을 뺀 코어 위치는 둘 다 \texttt{ERA}이고,
복원된 키로 설정한 에니그마는 실제 키의 에니그마와 똑같이 동작한다. 오른쪽 링 \texttt{C}는 정확히 맞았다.
가운데 링이 \texttt{M}이 아니라 \texttt{A}로 나온 것은 이 전문(121글자) 동안 왼쪽 로터가 한 번도 움직이지
않아서 가운데 링을 정할 정보가 전문에 없었기 때문이다.

\section{점수: 독일어다운 정도}

\begin{lstlisting}[language=Rust,caption={\code{check.rs}: 복호문의 점수.}]
pub fn score_german(text: &[u8]) -> f64 {
    if text.len() < 2 {
        return f64::NEG_INFINITY;
    }
    let t = bigram_log_probs();
    text.windows(2).map(|w| t[w[0] as usize][w[1] as usize]).sum::<f64>() / (text.len() - 1) as f64
}
\end{lstlisting}

\code{text.windows(2)}는 이웃한 두 글자의 쌍을 차례로 내놓는다. 쌍마다 ``앞 글자 다음에 뒤 글자가 올
확률''의 로그를 표에서 읽어 평균을 낸다. 표는 독일어 글에서 두 글자 쌍(\term{바이그램}{bigram})의 빈도를
세어 만든 것이다. 무작위 글자열은 약 $-3.5$, 독일어 문장은 약 $-2.4$가 나온다.

\Why 글자 하나하나의 빈도가 아니라 \emph{이웃한 두 글자}를 보는 데에는 이유가 있다. 처음에는 글자의 빈도만
보았는데, 그 점수로는 다음 단계에서 스테커를 복원하지 못했다. 글자의 빈도는 글자의 순서를 보지 않는다.
스테커 한 쌍이 틀려 몇 글자가 다른 글자로 바뀌어도 전체 빈도는 거의 그대로다. 두 글자 쌍을 보면 틀린 글자
하나가 앞뒤 두 쌍을 독일어에 없는 조합으로 만들므로 점수에 뚜렷이 드러난다.

\begin{lesson}{점수 함수는 풀려는 문제의 글과 닮아야 한다}
바이그램으로 바꾼 뒤에도 전문의 첫 세 글자가 틀리는 경우가 남았다. 올바른 \texttt{ANX}가 아니라 \texttt{SLE}나
\texttt{AGI}가 나왔다. 가운데 링의 후보들이 전문의 앞부분에서만 서로 다른 복호문을 주는데, 올바른 복호문에는
군용 전문이 낱말 사이에 넣는 \texttt{X}가 들어 있었다. 바이그램 표는 문학 작품에서 만든 것이어서
\texttt{NX}나 \texttt{XO} 같은 쌍을 거의 본 적이 없었고, 그래서 올바른 답에 나쁜 점수를 주었다.
같은 글에서 낱말 사이를 \texttt{X}로 바꿔 한 번 더 센 것을 표에 더하자 이 경우까지 맞게 되었다.
\end{lesson}

\section{4단계: 남은 스테커}

크립 구간의 평문에도 암호문에도 나오지 않은 글자는 1, 2단계의 논리로 스테커를 정할 수 없다. 이 글자들에
대해서는 \term{언덕 오르기}{hill climbing}를 한다. 정해지지 않은 글자 두 개를 짝지어 보고, 전문 전체의
점수가 분명히(0.01 넘게) 오르는 짝 가운데 가장 많이 오르는 것을 채택한다. 더 오르는 짝이 없을 때까지
되풀이한다.

\General 언덕 오르기는 안개 속에서 산꼭대기를 찾는 방법이다. 지금 선 자리에서 한 걸음씩 디뎌 보고 가장 많이
올라가는 쪽으로 간다. 어느 쪽으로도 올라가지 않으면 멈춘다. 전체를 다 살피지 않고도 좋은 답에 이르지만,
그것이 가장 높은 꼭대기라는 보장은 없다.

\Ours 이 단계는 추론이 아니라 추측이다. 그래서 결과에 따로 표시한다. \code{unresolved}는 논리로 정하지
못한 글자들이고 \code{guessed}는 그 가운데 추측으로 짝을 지은 글자들이다. 예제에서는 \texttt{X}와
\texttt{Y}가 정해지지 않았고, 점수가 오르는 짝이 없어 스테커가 없는 것으로 두었다. 실제 키에서도 그 두
글자에는 플러그가 꽂혀 있지 않았다.

\section{전체 탐색: \texorpdfstring{\code{search}}{search}}

\begin{lstlisting}[language=Rust,caption={\code{search.rs}: 일을 나눠 병렬로 돌린다(일부).}]
let results = par_map(&tasks, opts.parallel, |&(i, t, g)| {
    let bo = BombeOptions { diagonal_board: opts.diagonal_board, turnover_at: t };
    let res = run_slice(&plugged, &orders[i], &bo, g);
    let sols: Vec<Solution> =
        res.stops.iter().flat_map(|s| check_stop(ciphertext, menu, s, opts.diagonal_board)).collect();
    (res, sols)
});
\end{lstlisting}

\code{tasks}는 (로터 순서, 넘어감 가설, 그리스 로터 위치)의 모든 조합이다. \code{par\_map}은 이 일들을
컴퓨터의 모든 처리 장치 코어에 나눠 준다. 일 하나는 봄베를 돌리고(\code{run\_slice}) 나온 정지를 모두
체킹한다(\code{check\_stop}).

\Why 이렇게 나눌 수 있는 것은 일들이 서로 독립이기 때문이다. 한 로터 순서의 결과는 다른 로터 순서의 결과에
아무 영향을 주지 않는다. 공유하는 것은 읽기만 하는 메뉴와 암호문뿐이다.

\General 계산대가 여러 개인 가게와 같다. 손님(일)들이 서로 상관없으므로 계산대를 늘리는 만큼 빨라진다.

\Ours 블레츨리 파크에서도 같은 방식이었다. 봄베 여러 대가 서로 다른 로터 순서를 나눠 맡았다. 그리고 이것은
F부에서 볼 ``회로에서는 동시가 기본''(\ref{sec:concurrent}절)의 소프트웨어 쪽 모습이기도 하다. FPGA에서는 엔진 여러 대를 나란히
놓는 것으로 같은 일을 한다. 측정 결과, 논리 스레드 24개의 PC에서 Enigma~I의 120가지 로터 순서(약 211만
위치)는 0.2초대에, M4의 1,344가지(약 6억 1천만 위치)는 98초에 끝난다.
